% The numerical core (C++): containers, filters, forecasters, optimizers, risk, the portfolio engine,
% metrics and the regression engine. Excerpts and locations come from extract_snippets.py.

\chapter{Time series containers}\label{c:sec:dataframe}

All numerical code works on Eigen matrices, but financial data is indexed by time.
\code{TimeSeries<Index>} and \code{DataFrame<Index>} pair an Eigen vector or matrix (rows
are times) with a \emph{strictly increasing} index:
\begin{itemize}
  \item The index type satisfies the \code{TimeIndex} concept (copyable and totally
        ordered): calendar dates, timestamps or integers.
  \item Construction validates the shapes, the ordering and unique column names.
  \item Time lookups use binary search.
\end{itemize}

\paragraph{Alignment.} Markets have different holidays. A cross-asset backtest keeps only
the dates where every asset has a finite price.

\paragraph{Returns.} Simple returns $r_t=P_t/P_{t-1}-1$ (\code{pct\_change}) or log returns
$\log(P_t/P_{t-1})$ are computed between consecutive \emph{kept} rows. A return can
therefore span a removed holiday. That is the right convention for a portfolio that can
only trade when all its markets are open.

\snippet{df_drop_missing}{C++}{Alignment: keep rows where every column is finite}
\begin{codepointers}
  \codeitem{time_index}{\code{TimeIndex} concept}
  \codeitem{timeseries}{\code{TimeSeries}}
  \codeitem{dataframe}{\code{DataFrame}}
  \codeitem{pct_change}{simple returns}
  \codeitem{log_returns}{log returns}
\end{codepointers}

\chapter{Signal extraction}\label{c:sec:signal}

Filters implement \code{IFilter<P>}. They are \emph{causal}: output $t$ uses inputs up to
$t$ only. Two modes are supported:
\begin{itemize}
  \item streaming: \code{process(x)}, one sample at a time, keeping state;
  \item batch: \code{apply(signal)}, from a fresh state, without disturbing the streaming
        state.
\end{itemize}
Filter design is delegated to iir1, behind a private implementation struct (pimpl), so iir1
never appears in public headers.

\section{Exponentially weighted moving average}
The EWMA is the one-pole low-pass
\[
  y_t=\alpha x_t+(1-\alpha)y_{t-1},\qquad H(z)=\frac{\alpha}{1-(1-\alpha)z^{-1}},\qquad \alpha\in(0,1].
\]
Its DC gain is $H(1)=1$, so a constant input passes unchanged.

\begin{proposition}[Seeding]
Let $\tilde y$ be the EWMA of $x_t-x_0$ started from a zero state. Then $y_t=x_0+\tilde y_t$
equals the recursion started at $y_{-1}=x_0$:
\[
  y_t=\alpha\sum_{k=0}^{t}(1-\alpha)^k x_{t-k}+(1-\alpha)^{t+1}x_0 .
\]
\end{proposition}
\begin{proof}
Unrolling the recursion from $y_{-1}$ gives
$y_t=\alpha\sum_{k=0}^{t}(1-\alpha)^kx_{t-k}+(1-\alpha)^{t+1}y_{-1}$. With a zero state,
$\tilde y_t=\alpha\sum_{k=0}^t(1-\alpha)^k(x_{t-k}-x_0)$. Using
$\alpha\sum_{k=0}^t(1-\alpha)^k=1-(1-\alpha)^{t+1}$, adding $x_0$ gives the formula with
$y_{-1}=x_0$.
\end{proof}
Starting from zero instead would bias the early outputs towards zero, so the filter seeds
with the first sample by default.

\paragraph{Parameterizations.} Two helpers convert the usual parameterizations to $\alpha$:
\begin{itemize}
  \item \code{alpha\_from\_span}$(N)=2/(N+1)$. The mean age of the EWMA weights is
        $\sum_k k\,\alpha(1-\alpha)^k=(1-\alpha)/\alpha$, and that of an $N$-point moving
        average is $(N-1)/2$. Equating the two gives $\alpha=2/(N+1)$.
  \item \code{alpha\_from\_halflife}$(h)=1-2^{-1/h}$. This solves $(1-\alpha)^h=\tfrac12$:
        the weight $h$ steps back is half the newest weight.
\end{itemize}

\snippet{ewma_impl}{C++}{EWMA on iir1's one-pole section, with first-sample seeding}
\begin{codepointers}
  \codeitem{ewma_class}{class and parameters}
  \codeitem{ewma_alphas}{span and half-life conversions}
  \codeitem{moving_average}{simple moving average helper (rolling sum, $O(n)$)}
\end{codepointers}

\section{Butterworth band-pass}
A causal Butterworth band-pass isolates trend cycles between two frequencies
$0<f_\ell<f_h<f_s/2$. It is built in three steps.
\begin{enumerate}
  \item \textbf{Analog prototype.} The order-$N$ low-pass prototype is maximally flat:
        \[
          |H(j\Omega)|^2=\frac{1}{1+\Omega^{2N}},
        \]
        with poles $s_k=e^{j\pi(2k+N-1)/(2N)}$, $k=1,\dots,N$, on the left half of the unit
        circle.
  \item \textbf{Low-pass to band-pass.} The substitution
        \[
          s\mapsto\frac{s^2+\Omega_0^2}{sB}
        \]
        maps it to a band-pass with $2N$ poles, centred on $\Omega_0$ with bandwidth $B$.
  \item \textbf{Bilinear transform.} The map $s=2f_s\frac{1-z^{-1}}{1+z^{-1}}$, with the band
        edges prewarped by $\Omega=2f_s\tan(\pi f/f_s)$, gives the digital filter. It is run
        as a cascade of second-order sections, which is numerically stable at high orders.
\end{enumerate}
iir1 takes the band as a centre and a width, so the edges map to centre
$(f_\ell+f_h)/2$ and width $f_h-f_\ell$. The filter then has $|H|=1/\sqrt2$ ($-3$\,dB)
at both edges and unit gain at the centre. The tests check exactly these properties, and
strong attenuation outside the band. Being causal, the filter delays the signal: the
group delay grows with the order and with narrower bands.
\begin{codepointers}
  \codeitem{butterworth_params}{parameters (order, sample rate, edges)}
  \codeitem{butterworth_design}{design: edges to centre and width}
\end{codepointers}

\chapter{Forecasting}\label{c:sec:forecast}

Forecasters implement \code{IForecaster<P>}:
\begin{itemize}
  \item \code{fit(y)} estimates the model on a series;
  \item \code{forecast(H)} returns the mean $\hat y_{T+h}$ and the predictive variance $v_h$
        for $h=1..H$;
  \item Gaussian prediction intervals follow as $\hat y\pm z_{(1+c)/2}\sqrt{v_h}$.
\end{itemize}
Chapters~\ref{ch:es}--\ref{ch:gp} derive each model in full. This section states
what the code computes.

\section{Simple exponential smoothing}
The model is the ETS(A,N,N) state space
\[
  y_t=\ell_{t-1}+\varepsilon_t,\qquad \ell_t=\ell_{t-1}+\alpha\varepsilon_t,\qquad
  \varepsilon_t\sim\N(0,\sigma^2).
\]
The fit initializes $\ell=y_1$ and runs the error-correction form
$\ell\leftarrow\ell+\alpha(y_t-\ell)$, accumulating $\mathrm{SSE}=\sum_t(y_t-\ell_{t-1})^2$.
It then estimates $\hat\sigma^2=\mathrm{SSE}/T$.

\begin{proposition}
The $h$-step forecast is $\hat y_{T+h}=\ell_T$, with variance $\sigma^2\big(1+(h-1)\alpha^2\big)$.
\end{proposition}
\begin{proof}
Iterating the state equation gives $\ell_{T+h-1}=\ell_T+\alpha\sum_{j=1}^{h-1}\varepsilon_{T+j}$,
so $y_{T+h}=\ell_T+\alpha\sum_{j=1}^{h-1}\varepsilon_{T+j}+\varepsilon_{T+h}$. The errors are
independent, so the variance is $(h-1)\alpha^2\sigma^2+\sigma^2$.
\end{proof}
\snippet{ses_fit}{C++}{SES fit and forecast variance}

\section{Kalman filter}
The model is the linear Gaussian state space
\[
  x_t=Fx_{t-1}+w_t,\quad w_t\sim\N(0,Q);\qquad y_t=Hx_t+v_t,\quad v_t\sim\N(0,R),
\]
with prior $x_0\sim\N(\hat x_0,P_0)$, where the prior describes the first state. Each step
first predicts:
\[
  \hat x_{t|t-1}=F\hat x_{t-1|t-1},\qquad P_{t|t-1}=FP_{t-1|t-1}F^\T+Q .
\]
Then it computes the innovation and its variance:
\[
  e_t=y_t-H\hat x_{t|t-1},\qquad S_t=HP_{t|t-1}H^\T+R .
\]
Finally it corrects, with gain $K_t=P_{t|t-1}H^\T/S_t$:
\[
  \hat x_{t|t}=\hat x_{t|t-1}+K_te_t,\qquad
  P_{t|t}=(I-K_tH)P_{t|t-1}(I-K_tH)^\T+R\,K_tK_t^\T .
\]
\begin{itemize}
  \item \textbf{Joseph form.} The covariance update above is the posterior covariance for
        \emph{any} gain. It is symmetric and positive semi-definite by construction, so it
        stays valid under rounding, unlike the short form $(I-KH)P$.
  \item \textbf{Missing observations.} Non-finite $y_t$ skip the correction.
  \item \textbf{Log-likelihood.}
        $\ell=-\tfrac12\sum_t\big(\log2\pi+\log S_t+e_t^2/S_t\big)$, summed over the observed
        points.
  \item \textbf{Forecasts.} They iterate the prediction step: mean $Hx$, variance
        $HPH^\T+R$.
  \item \textbf{Smoothing.} The Rauch--Tung--Striebel smoother uses the gain
        $C_t=P_{t|t}F^\T P_{t+1|t}^{-1}$ (computed by an $LDL^\T$ solve, not an inverse) and
        $\hat x^s_t=\hat x_{t|t}+C_t(\hat x^s_{t+1}-\hat x_{t+1|t})$.
\end{itemize}

Two ready-made models are provided:
\begin{itemize}
  \item \emph{local level}: $F=H=1$, so the level is a random walk observed with noise;
  \item \emph{local linear trend}: $F=\begin{psmallmatrix}1&1\\0&1\end{psmallmatrix}$,
        $H=(1,0)$, $Q=\diag(q_\ell,q_s)$, so level and slope both drift.
\end{itemize}
Both use a diffuse prior ($P_0=10^7I$). The local linear trend's slope state is a direct
estimate of the trend.
\snippet{kalman_step}{C++}{Kalman predict and Joseph-form correction}
\begin{codepointers}
  \codeitem{kalman_params}{parameters $F,H,Q,R,x_0,P_0$}
  \codeitem{kalman_models}{local level and local linear trend}
  \codeitem{kalman_forecast}{forecast}
  \codeitem{kalman_smoother}{RTS smoother}
\end{codepointers}

\section{Gaussian process regression}
The series is regressed on time $t=0,\dots,T-1$ in three steps.
\begin{enumerate}
  \item \textbf{Standardization.} The observations are standardized,
        $z=(y-\bar y)/s$, where $s$ is the sample standard deviation when normalizing.
  \item \textbf{Model.} $z\sim\N(0,K)$ with $K=K_f+\sigma_n^2I$ and
        $(K_f)_{ij}=k(t_i,t_j)$. The available kernels, with $r=t-t'$:
\begin{align*}
  k_{\text{SE}}&=\sigma^2\exp\!\big(-r^2/2\ell^2\big),\\
  k_{\text{M52}}&=\sigma^2\big(1+u+u^2/3\big)e^{-u},\quad u=\sqrt5|r|/\ell,\\
  k_{\text{per}}&=\sigma^2\exp\!\big(-2\sin^2(\pi r/p)/\ell^2\big).
\end{align*}
        Sums and products of kernels are kernels too.
  \item \textbf{Hyperparameters.} With $\alpha=K^{-1}z$ (from a Cholesky factorization
        $K=LL^\T$), the log marginal likelihood is
\[
  \log p(z)=-\tfrac12z^\T\alpha-\textstyle\sum_i\log L_{ii}-\tfrac T2\log2\pi .
\]
        The hyperparameters $\theta$ (logs of variances, length scales, periods and,
        optionally, the noise) maximize it, using L-BFGS with the analytic gradient
\[
  \frac{\partial(-\log p)}{\partial\theta_j}=-\tfrac12\tr\!\Big(\big(\alpha\alpha^\T-K^{-1}\big)\frac{\partial K}{\partial\theta_j}\Big).
\]
        Working in logs keeps the parameters positive. Values with $|\theta_j|>20$, or a
        $K$ that is not positive definite, are returned as $+\infty$, which the line search
        rejects. The noise has a floor (\code{min\_noise\_variance}).
\end{enumerate}

\paragraph{Prediction.} At new times $t_*$, with $v=L^{-1}k_*$:
\[
  \text{mean}=\bar y+s\,k_*^\T\alpha,\qquad
  \text{variance}=s^2\big(k(t_*,t_*)-v^\T v+\sigma_n^2\big).
\]
The reported likelihood is converted back to the units of $y$ by subtracting $T\log s$.
\snippet{gp_nll}{C++}{Gaussian process: negative log-likelihood gradient}
\begin{codepointers}
  \codeitem{gp_factorize}{Cholesky and log marginal likelihood}
  \codeitem{gp_fit}{fit with hyperparameter optimization}
  \codeitem{gp_predict}{predictive mean and variance}
  \codeitem{kernels}{kernels and their log-parameter gradients}
\end{codepointers}

\chapter{Optimization}\label{c:sec:optimization}

Optimizers implement \code{IOptimizer<P,R>}: \code{optimize(initial\_guess)} returns the
solution, and \code{result()} returns a solver-specific result object (status, iterations,
residuals). Linear constraints are expressed in matrix form:
$A_{\text{eq}}x=b_{\text{eq}}$ and $A_{\text{in}}x\le b_{\text{in}}$. Box bounds
$l\le x\le u$ convert to inequality rows.
\begin{codepointers}
  \codeitem{ioptimizer}{optimizer interface}
  \codeitem{constraints}{linear constraints and their conversion}
\end{codepointers}

\section{L-BFGS}\label{c:sec:lbfgs}
L-BFGS minimizes a smooth $f$ with search direction $d=-H_kg_k$. Here $H_k$ is the BFGS
inverse-Hessian approximation built from the last $m$ pairs
$s_j=x_{j+1}-x_j$, $y_j=g_{j+1}-g_j$, with $\rho_j=1/s_j^\T y_j$. It is applied in $O(mn)$
by the two-loop recursion:
\begin{enumerate}
  \item $q\leftarrow g$. For $j$ from newest to oldest: $a_j=\rho_js_j^\T q$, then
        $q\leftarrow q-a_jy_j$.
  \item $r\leftarrow\gamma q$, with $\gamma=s^\T y/y^\T y$ from the newest pair. On the
        first iteration $\gamma=1/\max(1,\norm{g})$.
  \item For $j$ from oldest to newest: $\beta=\rho_jy_j^\T r$, then
        $r\leftarrow r+(a_j-\beta)s_j$. Finally $d=-r$.
\end{enumerate}
Safeguards keep the method well defined:
\begin{itemize}
  \item A pair is stored only if $s^\T y>10^{-12}\norm s\norm y$ (the curvature
        condition), which keeps $H$ positive definite.
  \item If $d$ is not a descent direction, the memory is cleared and the method restarts
        from steepest descent.
  \item Iteration stops when $\norm{g}_\infty\le\varepsilon$.
\end{itemize}

\paragraph{Line search.} Backtracking from $t=1$ halves $t$ until the step is accepted. A
non-finite $f$ counts as rejection, which is how barrier and domain constraints are
enforced. A step is accepted under either condition:
\begin{description}
  \item[Armijo] $f(x+td)\le f(x)+c_1t\,g^\T d$, with $c_1=10^{-4}$.
  \item[Approximate Wolfe] (Hager--Zhang), when $f$ is flat to within rounding:
        $f(x+td)\le f(x)+\epsilon_f|f(x)|$ and
        $\sigma\,g^\T d\le g(x+td)^\T d\le(2c_1-1)\,g^\T d$, with $\epsilon_f=10^{-10}$ and
        $\sigma=0.9$.
\end{description}
Near a minimum the decrease Armijo requires is smaller than the rounding error in $f$,
while directional derivatives remain accurate. The second test lets the method keep
converging there. On 500-asset risk budgeting it took the attainable budget accuracy from
about $3\cdot10^{-8}$ to $10^{-8}$.
\snippet{lbfgs_two_loop}{C++}{Two-loop recursion}
\snippet{lbfgs_line_search}{C++}{Line search: Armijo or approximate Wolfe}
\begin{codepointers}
  \codeitem{lbfgs_params}{parameters and defaults}
  \codeitem{lbfgs_optimize}{full algorithm}
\end{codepointers}

\section{Augmented Lagrangian}\label{c:sec:al}
For $\min f(x)$ subject to $h(x)=A_{\text{eq}}x-b_{\text{eq}}=0$ and
$g(x)=A_{\text{in}}x-b_{\text{in}}\le0$, the Powell--Hestenes--Rockafellar augmented
Lagrangian is
\[
  L_\rho(x;\lambda,\mu)=f(x)+\lambda^\T h+\frac\rho2\norm h^2
  +\frac1{2\rho}\sum_i\Big(\max(0,\mu_i+\rho g_i)^2-\mu_i^2\Big),
\]
with gradient
$\nabla f+A_{\text{eq}}^\T(\lambda+\rho h)+A_{\text{in}}^\T\max(0,\mu+\rho g)$.
Each outer iteration does four things:
\begin{enumerate}
  \item \textbf{Minimize} $L_\rho$ with L-BFGS, warm-started at the previous iterate. The
        inner tolerance starts at $\max(\varepsilon,0.1)$ and shrinks tenfold per outer
        iteration down to $\varepsilon$.
  \item \textbf{Measure} the violation
        $v=\max\big(\norm h_\infty,\norm{\max(g,-\mu/\rho)}_\infty\big)$. The second term is
        zero if and only if $g\le0$ and $\mu_ig_i=0$, so it also measures complementarity.
  \item \textbf{Update the multipliers}: $\lambda\leftarrow\lambda+\rho h$ and
        $\mu\leftarrow\max(0,\mu+\rho g)$. After this update,
        $\nabla_xL_\rho=\nabla f+A_{\text{eq}}^\T\lambda+A_{\text{in}}^\T\mu$, so the inner
        gradient norm is the KKT stationarity.
  \item \textbf{Test and adapt.} Stop when $v\le\varepsilon$ and stationarity
        $\le\varepsilon$. Otherwise, if $v>0.25\,v_{\text{prev}}$, set
        $\rho\leftarrow\min(10\rho,10^{10})$ (starting from $\rho_0=10$).
\end{enumerate}
\snippet{al_function}{C++}{Augmented Lagrangian value and gradient}
\snippet{al_updates}{C++}{Multiplier and penalty updates}
\begin{codepointers}
  \codeitem{al_options}{options and defaults}
  \codeitem{al_optimize}{full algorithm}
\end{codepointers}

\section{Sequential quadratic programming}
At $x_k$, SQP solves the quadratic subproblem
\[
  \min_d\;g_k^\T d+\tfrac12d^\T B_kd\quad\text{s.t.}\quad
  A_{\text{eq}}d=b_{\text{eq}}-A_{\text{eq}}x_k,\;\;A_{\text{in}}d\le b_{\text{in}}-A_{\text{in}}x_k .
\]
It uses the augmented Lagrangian of \S\ref{c:sec:al} at tolerance $10^{-10}$, or a Cholesky
solve $d=-B^{-1}g$ when there are no constraints.
\begin{description}
  \item[Convergence] is declared when three KKT residuals at $x_k$, computed with the QP
        multipliers, are all below $\varepsilon$:
        \begin{itemize}
          \item stationarity
                $\norm{g+A_{\text{eq}}^\T\lambda+A_{\text{in}}^\T\mu}_\infty$;
          \item constraint violation;
          \item complementarity $\max_i|\mu_ig_i|$.
        \end{itemize}
  \item[Merit function] $\phi(x)=f(x)+\nu V_1(x)$, with
        $V_1=\norm h_1+\sum_i\max(g_i,0)$. The penalty is kept at
        $\nu\ge\norm{(\lambda,\mu)}_\infty+10^{-3}$. Because the constraints are linear and
        $d$ satisfies the linearized constraints, $V_1(x+td)\le(1-t)V_1(x)$. The
        directional derivative is therefore bounded by $g^\T d-\nu V_1(x)<0$, so $d$
        descends $\phi$. Backtracking applies Armijo to $\phi$ with this slope.
  \item[Restart] If the line search fails, $B$ resets to the identity. A failure from
        $B=I$ ends the iteration.
\end{description}

\begin{proposition}[Damped BFGS, Powell]
Let $r=\theta y+(1-\theta)Bs$, with $\theta=1$ if $s^\T y\ge0.2\,s^\T Bs$ and
$\theta=0.8\,s^\T Bs/(s^\T Bs-s^\T y)$ otherwise. The update
\[
  B^+=B-\frac{Bss^\T B}{s^\T Bs}+\frac{rr^\T}{s^\T r}
\]
keeps $B$ positive definite.
\end{proposition}
\begin{proof}
If $\theta<1$, then $s^\T r=s^\T Bs-\theta(s^\T Bs-s^\T y)=0.2\,s^\T Bs>0$. If $\theta=1$,
then $s^\T r=s^\T y\ge0.2\,s^\T Bs>0$. The BFGS update with a pair satisfying $s^\T r>0$
preserves positive definiteness.
\end{proof}
After the first step with $s^\T y>0$, $B$ is scaled to $(y^\T y/s^\T y)I$ before the update,
which matches the curvature along $s$.
\snippet{sqp_bfgs}{C++}{Damped BFGS update}
\snippet{sqp_merit}{C++}{$\ell_1$ merit penalty and line search}
\begin{codepointers}
  \codeitem{sqp_options}{options and defaults}
  \codeitem{sqp_qp}{QP subproblem through the augmented Lagrangian}
  \codeitem{sqp_optimize}{full algorithm}
\end{codepointers}

\section{Nelder--Mead with box bounds}
Nelder--Mead keeps a simplex of $n+1$ vertices sorted by value,
$f(x_0)\le\dots\le f(x_n)$, with $c$ the centroid of the best $n$. Each iteration:
\begin{enumerate}
  \item \textbf{Reflect}: $x_r=c+\alpha(c-x_n)$.
  \item If $f_r<f_0$, \textbf{expand} to $x_e=c+\gamma(x_r-c)$ and keep the better of
        $x_e$ and $x_r$.
  \item Else, if $f_r<f_{n-1}$, accept $x_r$.
  \item Else \textbf{contract}:
        \begin{itemize}
          \item outside, $x_c=c+\rho(x_r-c)$, if $f_r<f_n$ (accepted if $f_c<f_r$);
          \item inside, $x_c=c+\rho(x_n-c)$, otherwise (accepted if $f_c<f_n$).
        \end{itemize}
  \item Else \textbf{shrink} every vertex towards the best:
        $x_j\leftarrow x_0+\sigma(x_j-x_0)$.
\end{enumerate}
The \emph{adaptive} coefficients of Gao and Han, $\alpha=1$, $\gamma=1+2/n$,
$\rho=\tfrac34-\tfrac1{2n}$, $\sigma=1-\tfrac1n$, avoid the stagnation of the standard
values $(1,2,\tfrac12,\tfrac12)$ in higher dimensions.

Bounds are enforced by projecting every trial point onto the box. The box is convex, so a
shrink stays inside it.

\paragraph{Termination.} The iteration stops when both conditions hold:
\[
  f_n-f_0\le\tau_f(1+|f_0|)\quad\text{and}\quad\max_j\norm{x_j-x_0}_\infty\le\tau_x(1+\norm{x_0}_\infty).
\]
A NaN objective counts as $+\infty$.
\snippet{nm_step}{C++}{One Nelder--Mead iteration}
\begin{codepointers}
  \codeitem{nm_coefficients}{adaptive coefficients}
  \codeitem{box}{box projection}
  \codeitem{nm_optimize}{full algorithm, initial simplex}
\end{codepointers}

\section{Differential evolution}
\paragraph{Population.} $N_p$ vectors (default $15n$, at least 4, or 6 for rand/2), drawn
uniformly in the box. An initial guess, if given, replaces member 0.

\paragraph{Mutation.} Each generation draws one $F\sim U[F_{\min},F_{\max}]$ (dithering),
then builds a mutant per member $i$ from distinct random members $r_1,\dots,r_5\ne i$:
\begin{align*}
  \text{rand/1:}\quad& v=x_{r_1}+F(x_{r_2}-x_{r_3}),\\
  \text{best/1:}\quad& v=x_{\text{best}}+F(x_{r_1}-x_{r_2}),\\
  \text{current-to-best/1:}\quad& v=x_i+F(x_{\text{best}}-x_i)+F(x_{r_1}-x_{r_2}),\\
  \text{rand/2:}\quad& v=x_{r_1}+F(x_{r_2}-x_{r_3})+F(x_{r_4}-x_{r_5}).
\end{align*}

\paragraph{Crossover.} Binomial: $u_{ij}=v_{ij}$ if $U<\mathrm{CR}$ or $j=j_{\text{rand}}$,
else $x_{ij}$, so at least one coordinate comes from the mutant. A coordinate outside the
box bounces back to the midpoint between the parent and the violated bound.

\paragraph{Selection.} Greedy: $u$ replaces $x_i$ if $f(u)\le f(x_i)$. Ties go to the trial,
which keeps the population moving on plateaus.

\paragraph{Termination.} The run stops when the population's objective values have
$\mathrm{sd}(f)\le\tau_a+\tau|\overline f|$.

\paragraph{Parallelism and reproducibility.} All trial vectors of a generation are built
sequentially from one seeded RNG \emph{before} evaluation. Only the evaluations run in
parallel, so the result is the same for every thread count. A \code{ParallelEvaluator<T,R>}
does the evaluations:
\begin{itemize}
  \item each input becomes a \code{std::packaged\_task} in a queue;
  \item $n$ worker threads consume the queue;
  \item the caller collects the results, in input order, through the tasks' futures.
\end{itemize}
\snippet{de_generation}{C++}{Mutation and crossover of one generation}
\snippet{parallel_evaluator}{C++}{ParallelEvaluator: job queue, consumers, futures}
\begin{codepointers}
  \codeitem{de_params}{parameters and defaults}
  \codeitem{de_optimize}{full algorithm}
\end{codepointers}

\chapter{Risk measures}\label{c:sec:risk}

For portfolio weights $x$ on assets with return covariance $\Sigma$ and mean $\mu$, write
$\sigma(x)=\sqrt{x^\T\Sigma x}$ and take the loss $L=-r^\T x$ with $r\sim\N(\mu,\Sigma)$. Then
$L\sim\N(-\mu^\T x,\sigma(x)^2)$. The three measures are:
\begin{align*}
  \text{Volatility:}\quad & R(x)=\sigma(x), & \nabla R &= \Sigma x/\sigma(x),\\
  \text{VaR}_\alpha:\quad & R(x)=-\mu^\T x+z_\alpha\,\sigma(x), & \nabla R&=-\mu+z_\alpha\Sigma x/\sigma(x),\\
  \text{ES}_\alpha:\quad & R(x)=-\mu^\T x+\frac{\varphi(z_\alpha)}{1-\alpha}\,\sigma(x), & \nabla R&=-\mu+\frac{\varphi(z_\alpha)}{1-\alpha}\Sigma x/\sigma(x),
\end{align*}
with $z_\alpha=\Phi^{-1}(\alpha)$.

\begin{proposition}[Gaussian expected shortfall]
If $L\sim\N(m,s^2)$, then
$\mathrm{ES}_\alpha=\E[L\mid L\ge\mathrm{VaR}_\alpha]=m+s\,\varphi(z_\alpha)/(1-\alpha)$.
\end{proposition}
\begin{proof}
Write $L=m+sZ$. The event is $Z\ge z_\alpha$, which has probability $1-\alpha$. Since
$\varphi'(z)=-z\varphi(z)$, we get $\int_{z_\alpha}^\infty z\varphi(z)\,dz=\varphi(z_\alpha)$,
so $\E[Z\mid Z\ge z_\alpha]=\varphi(z_\alpha)/(1-\alpha)$.
\end{proof}

\begin{proposition}[Euler decomposition]
Each measure is positively homogeneous of degree one, $R(cx)=cR(x)$ for $c>0$. Hence
$R(x)=\sum_ix_i\,\partial_iR(x)$, and the \emph{risk contributions}
$\mathrm{RC}_i=x_i\partial_iR$ add up exactly to the total risk.
\end{proposition}
\begin{proof}
Differentiate $R(cx)=cR(x)$ with respect to $c$ at $c=1$.
\end{proof}
At $\sigma=0$ the gradient of $\sigma$ is set to zero, which is a valid subgradient of the
seminorm. The quantile $\Phi^{-1}$ uses Acklam's rational approximation (relative error
below $1.2\cdot10^{-9}$), refined to full precision by one Halley step on $\Phi(x)-p$:
\[
  x\leftarrow x-\frac{u}{1+xu/2},\qquad u=(\Phi(x)-p)\sqrt{2\pi}\,e^{x^2/2}.
\]
\snippet{vol_gradient}{C++}{Volatility gradient, the building block of every measure}
\snippet{var_measure}{C++}{Gaussian VaR and its gradient}
\begin{codepointers}
  \codeitem{irisk}{risk measure interface}
  \codeitem{risk_contributions}{risk contributions $x\odot\nabla R$}
  \codeitem{volatility_measure}{volatility}
  \codeitem{es_measure}{expected shortfall (constant $\varphi(z_\alpha)/(1-\alpha)$)}
  \codeitem{normal_quantile}{normal quantile: Halley refinement}
\end{codepointers}

\chapter{Risk budgeting}\label{c:sec:riskbudgeting}

Given budgets $b_i>0$ with $\sum_ib_i=1$, the \emph{risk budgeting portfolio} is the
long-only $x$ with $\mathrm{RC}_i(x)/R(x)=b_i$ for all $i$. Equal budgets give the equal risk
contribution portfolio.

\begin{theorem}\label{c:thm:rb}
Let $R$ be convex, positively homogeneous of degree one, differentiable on $y>0$, and
positive on $\{y\ge0,\,y\ne0\}$. Then
\[
  f(y)=R(y)-\sum_ib_i\log y_i
\]
has a unique minimizer $y^\ast>0$, and $x^\ast=y^\ast/\sum_iy^\ast_i$ is the unique long-only
portfolio with $\mathrm{RC}_i(x^\ast)/R(x^\ast)=b_i$. Moreover $R(y^\ast)=1$.
\end{theorem}
\begin{proof}
\emph{Existence.} $R$ is continuous and positive on the compact simplex
$\Delta=\{y\ge0,\sum y_i=1\}$, so $c=\min_\Delta R>0$. By homogeneity,
$R(y)\ge c\sum_iy_i$ for $y\ge0$. Hence $f(y)\ge\sum_i(cy_i-b_i\log y_i)$. Each term tends
to $+\infty$ as $y_i\to0^+$ or $y_i\to\infty$, so $f$ is coercive on $y>0$ and attains
its minimum.

\emph{Uniqueness.} $-\sum b_i\log y_i$ is strictly convex and $R$ is convex, so $f$ is
strictly convex.

\emph{Characterization.} At the minimizer, $\partial_if=\partial_iR(y^\ast)-b_i/y^\ast_i=0$,
so $\mathrm{RC}_i(y^\ast)=y^\ast_i\partial_iR(y^\ast)=b_i$. Summing and using Euler gives
$R(y^\ast)=\sum_ib_i=1$. Let $S=\sum_iy^\ast_i$. Because $\partial_iR$ is homogeneous of
degree zero, $x^\ast=y^\ast/S$ has $\mathrm{RC}_i(x^\ast)=b_i/S$ and $R(x^\ast)=1/S$, so
$\mathrm{RC}_i/R=b_i$.

Conversely, if $x$ satisfies the budget condition, then $y=x/R(x)$ satisfies
$\partial_if(y)=0$ and is therefore the unique minimizer.
\end{proof}

For VaR and ES, the condition ``$R>0$ on long-only portfolios'' needs
$z\,\sigma(x)>\mu^\T x$. The solver checks it at the start and throws otherwise.

\paragraph{Algorithm.} The steps follow \coderef{rb_solver}.
\begin{enumerate}
  \item \textbf{Normalize} the budgets. Empty budgets mean equal risk.
  \item \textbf{Start} at $y^0_i\propto b_i/\sigma_i$, scaled so that $R(y^0)=1$. That
        scale is the magnitude of the solution by the theorem. The start is the
        inverse-volatility portfolio for equal budgets, which is exact for uncorrelated
        assets.
  \item \textbf{Minimize} $f$ with L-BFGS (memory 20). Points with any $y_i\le0$ return
        $+\infty$, so the log barrier needs no explicit constraint.
  \item \textbf{Match the tolerance to the budget error.} L-BFGS stops on
        $\norm{\nabla f}_\infty$, but the target is the budget error
        $e=\max_i|\mathrm{RC}_i/R-b_i|$. Near the solution $R\approx1$ and
        $\mathrm{RC}_i-b_i=y_i\nabla_if$, so $e\lesssim\max_iy_i\,\norm{\nabla f}_\infty$.
        The inner tolerance therefore starts at $\tau/\max_iy^0_i$. If $e>\tau$ afterwards,
        L-BFGS warm-restarts with tolerance $\tau_{\text{in}}\cdot\tfrac12\tau/e$, at most five
        times.
  \item \textbf{Return} $x=y/\sum y$, with risk, contributions, budget error and a
        convergence flag.
\end{enumerate}
\snippet{rb_objective}{C++}{Objective and gradient, for any risk measure}
\snippet{rb_solve}{C++}{Tolerance matching on the budget error}
\begin{codepointers}
  \codeitem{rb_budgets}{budget normalization}
  \codeitem{rb_start}{starting point}
  \codeitem{rb_template}{\code{RiskBudgeting<Measure>} wrapper}
\end{codepointers}

\chapter{The walk-forward portfolio engine}\label{c:sec:engine}

\code{RiskBudgetingPortfolio} backtests a long-only risk budgeting portfolio from a
$T\times N$ matrix of simple returns, where row $t$ holds the returns over $(t-1,t]$.
Table~\ref{c:tab:portfolio} lists its parameters, which are also the JSON keys of the CLI and
the Python API.

\begin{table}[ht]\centering\small
\caption{Portfolio parameters: engine defaults and the defaults of the AI configuration.}
\label{c:tab:portfolio}
\begin{tabularx}{\linewidth}{@{}l l l X@{}}
\toprule
Key & Engine & AI & Meaning \\
\midrule
\code{covariance\_window} & 60 & 120 & rows $W$ per covariance estimate \\
\code{covariance\_estimator} & sample & exponential & equal or exponential weights \\
\code{ewma\_halflife} & 20 & 40 & half-life $h$ in rows (exponential) \\
\code{budgets} & equal & set by policy & risk budgets per asset \\
\code{risk\_measure} & volatility & volatility & volatility, \code{value\_at\_risk}, \code{expected\_shortfall} \\
\code{confidence\_level} & 0.99 & -- & VaR/ES level in $(0.5,1)$ \\
\code{use\_expected\_returns} & false & -- & VaR/ES: window mean instead of 0 \\
\code{target\_volatility} & none & 0.10 & annualized target $\sigma^\ast$ \\
\code{max\_leverage} & 2.0 & 3.0 & cap $L_{\max}$ on $\sum_iw_i$ \\
\code{rebalance\_frequency} & 1 & 21 & trade every $k$ rows \\
\code{transaction\_cost} & 0 & 0.0005 & cost $\kappa$ per unit traded (5\,bp) \\
\code{drift\_weights} & true & true & weights drift between trades \\
\code{initial\_capital} & 1 & $10^6$ & starting equity \\
\code{periods\_per\_year} & 252 & 252 & $P$, for annualization \\
\code{solver\_tolerance} & $10^{-8}$ & -- & budget error target $\tau$ \\
\code{solver\_max\_iter} & 5000 & -- & L-BFGS iteration budget \\
\code{num\_threads} & 0 (all) & -- & threads for the solves \\
\bottomrule
\end{tabularx}
\end{table}
\begin{codepointers}
  \codeitem{portfolio_params}{parameter struct}
  \codeitem{portfolio_validate}{validation}
  \codeitem{portfolio_engine}{the engine}
\end{codepointers}

\section{Timing and absence of look-ahead}
Rebalances happen at the close of rows $t_k=W-1+kf$ for $k=0,\dots,\lfloor(T-W)/f\rfloor$.
At rebalance $t$ the engine:
\begin{enumerate}
  \item estimates moments from rows $t-W+1,\dots,t$, all known at the close of $t$;
  \item solves for the risk budgeting weights and scales them to the volatility target;
  \item trades from the drifted weights to the targets.
\end{enumerate}
The new weights earn the returns of rows $t+1$ onwards. A test confirms there is no
look-ahead: overwriting the data from row 200 on leaves every weight up to row 200
bit-identical.

\section{Covariance estimation}
For the window's returns $r_1,\dots,r_W$ (oldest first) and weights $w_k\ge0$ with
$\sum w_k=1$:
\[
  \hat\mu=\sum_kw_kr_k,\qquad
  \hat\Sigma=\frac{\sum_kw_k(r_k-\hat\mu)(r_k-\hat\mu)^\T}{1-\sum_kw_k^2}.
\]
The weights are either equal ($w_k=1/W$) or exponential: $w_k\propto\lambda^{W-k}$ with
$\lambda=2^{-1/h}$, so the newest return has age 0 and weight 1 before normalization.

\begin{proposition}
For i.i.d.\ returns with covariance $\Sigma$, $\E\hat\Sigma=\Sigma$. With equal weights the
denominator is $(W-1)/W$, which recovers the usual $1/(W-1)$ estimator.
\end{proposition}
\begin{proof}
Write $r_k-\hat\mu=\sum_j(\delta_{kj}-w_j)(r_j-\mu)$. By independence,
$\E\big[(r_k-\hat\mu)(r_k-\hat\mu)^\T\big]=\sum_j(\delta_{kj}-w_j)^2\Sigma=(1-2w_k+\sum_jw_j^2)\Sigma$.
Weighting by $w_k$ and summing gives $(1-\sum_jw_j^2)\Sigma$.
\end{proof}
The quantity $1/\sum_kw_k^2$ is the effective number of observations.
\snippet{cov_estimate}{C++}{Window moments}

\section{Allocation, budget schedule and volatility targeting}
\paragraph{Budget schedule.} The budgets in force at rebalance row $t$ are the last
\emph{budget change} with first row at or before $t$. A schedule row dated $d$ takes effect
from the first returns row dated on or after $d$ (a binary search). This is how
time-varying budgets (from a rule, or from the sentiment pipeline of Part~\ref{part:sentiment}) enter one continuous backtest. Without a
schedule, the parameter budgets apply.

\paragraph{Volatility targeting.} The risk budgeting weights $x$ (summing to 1) are scaled
by the leverage
\[
  L=\min\!\Big(L_{\max},\;\frac{\sigma^\ast}{\sqrt{x^\T\hat\Sigma x\cdot P}}\Big),
\]
so the predicted annualized volatility is $\sigma^\ast$ unless the cap binds. Any weight
not invested, $1-\sum_iw_i$, is cash earning zero.
\snippet{schedule_rows}{C++}{Schedule dates to first rows}
\snippet{budgets_at}{C++}{Budgets in force at a rebalance}
\snippet{solve_target}{C++}{Target weights of one rebalance date}
\begin{codepointers}
  \codeitem{portfolio_allocate}{risk measure selection}
\end{codepointers}

\section{Phase 1: parallel solves}
The targets of all rebalance dates depend only on their own windows, so they are solved
first, in parallel. Each solve warm-starts from the previous date's unlevered weights,
which cuts iterations sharply because consecutive windows differ by a few rows. Warm
starts make solves depend on each other, so dates are grouped into \emph{fixed} blocks of 8
consecutive rebalances:
\begin{itemize}
  \item one block is one task;
  \item warm starts happen only within a block;
  \item the block boundaries do not depend on the thread count.
\end{itemize}
Results are therefore bit-identical for any number of threads. With 50 assets, daily
rebalancing and 10 years of data (2,271 solves), a backtest takes 375\,ms on one thread and
53\,ms on eight.
\snippet{engine_phase1}{C++}{Phase 1}
\snippet{parallel_for}{C++}{Thread pool helper}
\begin{codepointers}
  \codeitem{resolve_threads}{thread count (0 = hardware concurrency, capped by tasks)}
\end{codepointers}

\section{Phase 2: drift, trading and accounting}
Phase 2 walks through time. With held weights $w_t$ and asset returns $r_t$, each row does
the following:
\begin{align*}
  g_t&=w_t^\T r_t &&\text{gross return,}\\
  w&\leftarrow w_t\odot(1+r_t)/(1+g_t) &&\text{drift,}\\
  \tau_t&=\norm{w^{\text{target}}-w}_1,\quad c_t=\kappa\,\tau_t &&\text{turnover and cost (rebalance rows),}\\
  n_t&=(1+g_t)(1-c_t)-1,\quad E_t=E_{t-1}(1+n_t) &&\text{net return and equity.}
\end{align*}
\paragraph{Drift.} Holding $i$ is worth $w_iE$ and grows to $w_iE(1+r_i)$. Cash earns
nothing, so equity grows to $E(1+g)$. The new weight is the ratio of the two.

\paragraph{Costs.} They are paid out of equity at the close, after the period's return,
which gives $E(1+g)(1-c)$. The first trade buys the whole portfolio from cash, so its
turnover is $\sum_iw_i$.

\paragraph{Checks.} The engine stops with an error if $1+g_t\le0$ (the portfolio lost all
its capital). These identities are checked to $10^{-15}$ in the tests.
\snippet{engine_phase2}{C++}{Phase 2}

\chapter{Performance metrics}\label{c:sec:metrics}

The metrics use the net returns $n_1,\dots,n_T$ from the first trade on, with $P$ periods
per year and equity $E$ starting at 1.
\begin{center}\small
\begin{tabular}{@{}ll@{}}
\toprule
Metric & Definition \\
\midrule
Total return & $\prod_t(1+n_t)-1$ \\
Annualized return (CAGR) & $(1+\text{total})^{P/T}-1$ \\
Annualized volatility & $\mathrm{sd}(n)\sqrt P$ (sample standard deviation) \\
Sharpe ratio & $\bar n\,P/(\mathrm{sd}(n)\sqrt P)$ (zero risk-free rate) \\
Sortino ratio & $\bar n\,P/\big(\sqrt{\tfrac1T\sum_t\min(n_t,0)^2}\,\sqrt P\big)$ \\
Maximum drawdown & $\max_t\big(1-E_t/\max_{s\le t}E_s\big)$ \\
Max drawdown duration & longest run of periods below a previous peak \\
Calmar ratio & annualized return / maximum drawdown \\
Hit rate & fraction of periods with $n_t>0$ \\
Best / worst period & $\max_tn_t$, $\min_tn_t$ \\
Turnover & average $\tau$ per rebalance; annualized $\sum_t\tau_t\,P/T$ \\
Total costs & $\sum_tc_t$ \\
\bottomrule
\end{tabular}
\end{center}
\snippet{performance}{C++}{Performance metrics}
\begin{codepointers}
  \codeitem{drawdown_series}{drawdown series}
\end{codepointers}

\chapter{Elastic net factor regression}\label{c:sec:regression}

The regression module regresses a target series (a fund, a strategy, an asset) on candidate
factors. It chooses the factors, the penalty and the mix, and measures their uncertainty.
Typical uses are replicating a fund with tradable factors, explaining a strategy's returns,
and hedging. Chapter~\ref{ch:en} has every
proof. This section states the results and maps them to the code.

\section{The elastic net}
For a design $X\in\R^{n\times p}$ and a target $y$, with $\ell_1=\lambda\alpha$ and
$\ell_2=\lambda(1-\alpha)$, the solver minimizes the glmnet objective
\[
  \frac1{2n}\norm{y-b_0\mathbf 1-Xb}^2+\ell_1\norm b_1+\frac{\ell_2}2\norm b_2^2 .
\]
$\alpha=1$ is the lasso, $\alpha=0$ ridge, and $\lambda=0$ least squares.
\begin{itemize}
  \item \textbf{Intercept.} It is unpenalized, and its optimum is $b_0=\bar y-\bar x^\T b$.
        The solver works on centred data.
  \item \textbf{Standardization.} Each centred column is divided by its population
        standard deviation $s_j$, and the coefficients are reported as $b_j=\beta_j/s_j$.
        The penalty therefore charges a factor in proportion to its volatility, so units and
        leverage do not change the selection. Constant columns get $b_j=0$.
\end{itemize}

\begin{proposition}[Coordinate update]
With working columns $z_j$, $v_j=\norm{z_j}^2/n$ and residual $r$, the exact minimizer over
$\beta_j$ with the others fixed is
\[
  \beta_j\leftarrow\frac{\sth(\rho_j,\ell_1)}{v_j+\ell_2},\qquad
  \rho_j=\tfrac1nz_j^\T r+v_j\beta_j,\qquad
  \sth(z,t)=\operatorname{sign}(z)\max(|z|-t,0).
\]
\end{proposition}
\begin{proof}
In $\beta_j$ alone the objective is $\frac{v_j+\ell_2}2\beta_j^2-\rho_j\beta_j+\ell_1|\beta_j|$
plus a constant. Its subgradient condition $0\in(v_j+\ell_2)\beta_j-\rho_j+\ell_1\partial|\beta_j|$
has the stated solution in each of the cases $\rho_j>\ell_1$, $\rho_j<-\ell_1$ and
$|\rho_j|\le\ell_1$.
\end{proof}
Each update costs $O(n)$: the residual is corrected by $-z_j\Delta_j$. Cyclic coordinate
descent converges because the non-smooth term is separable (Tseng, 2001).
\begin{itemize}
  \item \textbf{Active set.} After a full sweep, the solver sweeps only the non-zero
        coefficients until they settle, then checks with a full sweep.
  \item \textbf{Stopping rule.} The solver stops when a full sweep moves no fitted value by
        more than a fraction $\tau$ of the target's variance:
        $\max_jv_j\Delta_j^2<\tau\norm{\tilde y}^2/n$.
\end{itemize}

\begin{proposition}[$\lambda_{\max}$]
For $\alpha>0$, $\beta=0$ is optimal if and only if
$\lambda\ge\lambda_{\max}=\max_j|z_j^\T\tilde y|/(n\alpha)$.
\end{proposition}
\begin{proof}
At $\beta=0$ the residual is $\tilde y$, and the optimality condition of an inactive
coefficient is $|z_j^\T r|/n\le\ell_1$.
\end{proof}
\begin{itemize}
  \item \textbf{Floor on $\alpha$.} As in glmnet, $\alpha$ is floored at $10^{-3}$, so ridge
        also gets a finite path.
  \item \textbf{Rounding.} The value is multiplied by $1+10^{-12}$, so rounding cannot
        leave a $10^{-17}$ coefficient at $\lambda_{\max}$.
  \item \textbf{Grid.} The automatic grid is $\lambda_i=\lambda_{\max}r^{i/(L-1)}$. It is
        solved from the largest $\lambda$ down, each fit warm-started from the previous
        one.
\end{itemize}

\paragraph{Degrees of freedom.} With active set $A$ and locally constant signs, the active
coefficients are an affine function of $y$:
$\beta_A=(Z_A^\T Z_A+n\ell_2I)^{-1}(Z_A^\T\tilde y-n\ell_1s_A)$. By Stein's lemma, the degrees of
freedom are the trace of the hat matrix:
\[
  \mathrm{df}=1+\tr\!\big(Z_A(Z_A^\T Z_A+n\ell_2I)^{-1}Z_A^\T\big)=1+\sum_i\frac{d_i}{d_i+n\ell_2},
\]
where the $d_i$ are the eigenvalues of $Z_A^\T Z_A$ and the 1 counts the intercept. The
lasso gives $|A|+1$ (Zou, Hastie and Tibshirani, 2007).

\paragraph{Grouping.} For standardized factors $i,j$ with correlation $\rho$ and coefficients
of the same sign, the optimality conditions give
$|\beta_i-\beta_j|\le\norm{\tilde y}\sqrt{2(1-\rho)}/(\sqrt n\,\ell_2)$. Near-duplicate factors
share a loading instead of one of them taking it arbitrarily. This is why $\alpha<1$ helps
with substitutes such as the S\&P 500 and the Nasdaq 100.
\snippet{enet_sweep}{C++}{One coordinate descent sweep}
\snippet{enet_lambda_max}{C++}{$\lambda_{\max}$ with the same arithmetic as the first update}
\begin{codepointers}
  \codeitem{enet_params}{parameters}
  \codeitem{enet_standardize}{centring and scaling}
  \codeitem{enet_cd}{solver loop: full and active-set sweeps}
  \codeitem{enet_path}{path with warm starts}
  \codeitem{enet_df}{effective degrees of freedom}
\end{codepointers}

\section{The regression engine}
\code{RegressionEngine} runs the following steps (\coderef{regression_run}):
\begin{enumerate}
  \item \textbf{Split} the $T$ rows in time order. The first
        $n=\lfloor fT+10^{-9}\rfloor$ rows train, and the rest test. The test set is an
        extrapolation. With $f=1$ the engine is in sample.
  \item For each candidate subset $S$ of $k_{\min}..k_{\max}$ factors:
        \begin{enumerate}
          \item choose $(\alpha,\lambda)$ by time-series cross-validation on the training
                rows;
          \item fit on all training rows, warm-started down the grid;
          \item score by the test MSE, or in sample by
                $\mathrm{BIC}=n\log\mathrm{MSE}+\mathrm{df}\log n$ or
                $\mathrm{AIC}=n\log\mathrm{MSE}+2\,\mathrm{df}$, which are $-2\log\hat L$
                under Gaussian errors plus the penalty, with constants dropped.
        \end{enumerate}
  \item \textbf{Keep} the lowest score. Optionally refit on all rows, and bootstrap.
\end{enumerate}

\paragraph{Cross-validation.} The training rows are cut into contiguous blocks with
boundaries $\lfloor kn/B\rfloor$, and $g$ rows are left out on each side of a validation
block.
\begin{itemize}
  \item \emph{Blocked $K$-fold} ($B=K$) trains on the rows before and after the validation
        block.
  \item \emph{Walk-forward} ($B=K+1$) trains only on the rows before it.
\end{itemize}
The estimate is the mean fold MSE, and its standard error is
$\sqrt{\widehat{\Var}_k(m_k)/K}$. The cell with the lowest mean wins. The one-standard-error
rule then moves, within that $\alpha$, to the largest $\lambda$ whose error is within one
standard error of the minimum, which gives a sparser model that is statistically as good.

\paragraph{Subset search.}
\begin{itemize}
  \item \emph{Exhaustive} evaluates $\sum_k\binom Nk$ subsets (1\,470 for 14 factors and
        sizes 1--4).
  \item \emph{Forward} adds, greedily, the factor that most improves the score, which takes
        $\sum_{k\le k_{\max}}(N-k+1)$ evaluations.
  \item \emph{Auto} chooses Exhaustive below \code{max\_exhaustive\_subsets}.
\end{itemize}
The test MSE of the winner of many subsets is optimistic, because a factor can win by noise.
The bootstrap below exposes such factors, and a later holdout gives a clean figure.

\paragraph{Moving block bootstrap.} Each of $B$ samples:
\begin{enumerate}
  \item concatenates training-row blocks of length $\ell$ with uniform starts until there
        are $n$ rows;
  \item refits the selected subset at the selected $(\alpha,\lambda)$;
  \item records the coefficients, the intercept and the test MSE.
\end{enumerate}
Blocks keep the serial dependence of returns, such as volatility clustering (K\"unsch, 1989).
The summaries are:
\begin{itemize}
  \item the mean and standard error;
  \item percentile intervals (linear interpolation at $q(B-1)$);
  \item selection frequencies, the share of samples with a non-zero coefficient;
  \item the spread of the test MSE.
\end{itemize}
The intervals are conditional on the selected model, so they do not include selection
uncertainty.

\paragraph{Determinism.} Subsets and bootstrap samples run on the thread pool of
\S\ref{c:sec:engine}. Each task writes its own slot, ties go to the first index, and sample $b$
seeds its own \code{mt19937\_64} from \code{seed\_seq(seed, b)}. Results are bit-identical
for any thread count.
\snippet{regression_folds}{C++}{Contiguous validation folds with a gap}
\snippet{regression_cv}{C++}{Cross-validation error, standard error and the 1-SE rule}
\snippet{regression_bootstrap_draw}{C++}{One moving block bootstrap sample}
\begin{codepointers}
  \codeitem{regression_params}{parameters (defaults and meaning)}
  \codeitem{regression_class}{engine class, index-free and \code{DataFrame} interfaces}
  \codeitem{regression_evaluate}{evaluation of one subset}
  \codeitem{regression_criterion}{BIC, AIC, MSE}
  \codeitem{regression_exhaustive}{exhaustive search}
  \codeitem{regression_forward}{forward search}
  \codeitem{regression_bootstrap}{bootstrap and its summaries}
\end{codepointers}
